\section{Limitations}

This study has several important limitations. First, MLL23 is the only internal development domain, and AML-LMU is the only external validation domain. A stronger external-validity claim would require multiple independent sources with harmonized metadata. Second, reliable patient or acquisition-group identifiers are not available in the working MLL23 manifests, and AML-LMU patient identifiers are not documented in the audited fields used here. The near-duplicate-aware split reduces a specific leakage risk but does not prove patient-level independence.

Third, the unknown classes are defined by the protocol. They are outside the known training taxonomy, not necessarily clinically abnormal or diagnostically positive. Fourth, class imbalance is substantial in both internal and external datasets, making macro-F1, balanced accuracy, per-class metrics, and AUPR orientation important. Fifth, the multiseed analysis uses five training seeds and one ResNet18 backbone family. This is stronger than a single-run report but not exhaustive over architectures, optimizers, augmentations, or loss configurations.

Sixth, ArcFace is evaluated with a single frozen scale and margin. The study intentionally does not tune ArcFace after the confirmatory result because doing so would reopen method development. Seventh, the MLL23 unknown set was inspected during earlier deliveries, so later internal analyses are not independent discovery experiments. Eighth, the AML-LMU mapping is defensible but still a protocol mapping between source abbreviations and the known taxonomy. Ninth, no external adaptation is performed, which strengthens the test-only claim but may underestimate what could be achieved by explicitly domain-adaptive methods.

Finally, persistent homology is not exhausted as a mathematical tool. The negative cubical PH result applies to the evaluated filtrations, vectorizers, and fusion designs. The Vietoris-Rips analysis is explanatory and exploratory, with weak and inconsistent topology--OSR correlations. It should not be interpreted as a validated topological predictor of open-set reliability.
